% 2026-11-16 站主問：「At the same time, judicial independence is not without its drawbacks？」
% Law 786--790 頁主對話自文字層讀過（787 頁原句、788--790 頁原句已核）。註 205--207 所引文獻未另取原件。
\chapter{司法獨立之代價}
\label{ch:70_independence}

\begin{qaquestion}
Law\ix{persons}{law@Law, David S.} 787 頁第五部的賽局以「At the same time, judicial independence is not without its drawbacks」〔同時，司法獨立並非沒有缺點〕起頭。
\begin{enumerate}
\tightlist
\item 這句話的前後文為何？Law 所說的缺點是什麼？
\item 他如何在賽局中處理這項缺點，又如何權衡？
\item 這一權衡有何不足？
\end{enumerate}
\end{qaquestion}

Law 所說的缺點，是法院本身也成為人民的代理人，帶來另一組代理成本。他在賽局中把這項成本化約為獨立的法院較常判政府違憲，誤判違憲因此增加；人民厭惡暴政，寧可承受誤判違憲。

\section{前後文}
\label{sec:70_independence-1}

786 頁的結論是：只有獨立的法院所提供的資訊可信，人民若知道槍手拿槍指著法官的頭或挾持其家人，便不會重視法官的報告，因此人民有充分理由支持司法獨立\ix{terms}{02@\zhlidxcat{司法權之基礎}!司法獨立}。787 頁緊接著轉折：「To introduce judicial review by an independent court is, in effect, to trade one set of agency costs for another」〔引進由獨立法院行使的司法審查，實際上是以一組代理成本\ix{terms}{01@\zhlidxcat{賽局理論}!代理成本}換另一組〕。人民與政府之間的委託代理關係本有嚴重的代理成本，司法審查是處理這些成本的手段之一；但人民因此建立另一個代理關係，「no less than the government that it is supposed to monitor, the court itself is an agent that may attempt to advance its own goals at the expense of its principal」〔與它所監督的政府一樣，法院本身也是代理人，可能犧牲委託人而追求自己的目標〕。註 207 提及監督者自身的代理成本（agency costs in the monitors）。\footcite[786--787 \& n.207]{law2009}

\section{賽局中的處理與權衡}
\label{sec:70_independence-2}

\begin{enumerate}
\tightlist
\item
  賽局設定：人民選擇是否使法院獨立，政府選擇合憲或違憲的作為，法院判政府勝或敗；人民選擇司法獨立，即提高法院判政府敗訴的機率。人民不能自行判斷政府行為是否合憲，也無從得知法院的訊號是否正確；政府與法院則有完全資訊。\footcite[787--788]{law2009}
\item
  結果分四種：政府合憲而法院判合憲、政府違憲而法院判合憲（誤判合憲）、政府合憲而法院判違憲（誤判違憲）、政府違憲而法院判違憲。司法獨立提高的是判違憲的比率，並未提高判對的比率：「If the public opts for judicial independence, the incidence of false negatives will fall, but the incidence of false positives will rise」〔人民若選擇司法獨立，誤判合憲將減少，誤判違憲將增加〕；法院若受制於政府，情形相反。人民只能在兩害之間取其輕。\footcite[788--789]{law2009}
\item
  權衡：誤判合憲讓違憲的政府逍遙，可能釀成暴政；過度積極的法院造成的成本「more easily endured」〔較易承受〕，理由是「Unlike the government, a constitutional court has no obvious means of oppressing the people on its own」〔與政府不同，憲法法院沒有自行壓迫人民的明顯手段〕。\footcite[789--790]{law2009}由此推出：人民即使在法院推翻受歡迎的政策時，仍會支持司法獨立；只要相信法院比人民更懂政府行為是否合憲，或者把一定程度不想要的司法造法當成換取司法監督所須付的代價，這樣做就是理性的。\footcite[790]{law2009}
\end{enumerate}

\section{權衡的不足}
\label{sec:70_independence-3}

\begin{enumerate}
\tightlist
\item
  787 頁所說的代理成本，是法院追求自己的目標；賽局中的代理成本，只剩下判違憲的比率偏高。法院追求的目標未必是多判政府敗訴：法官的意識形態、任命者的期待、法院自身權力的擴張，可能使法院在某類案件多判違憲，在另一類案件多判合憲，例如註 238 所舉總統偏好任命主張擴張總統權的法官\pg{798 頁}。\footcite[798 n.238]{law2009}第 3 部所論法官選任，正是處理這一項代理成本的制度，賽局本身沒有納入。
\item
  「沒有壓迫手段」的論據，把法院的成本限於壓迫。法院宣告法律違憲，即排除民選機關所選擇的政策；Law 自己在 790 頁把這稱為不想要的司法造法。這項成本由希望該政策施行的那一部分人民承擔，屬於分配上的成本；人民分裂時（第 1 部第三節第 2 項），對這部分人民而言，誤判違憲未必較易承受。
\item
  人民無從判斷法院的訊號是否正確\pg{788 頁}，也就無從得知法院何時在追求自己的目標。Law 承認賽局未考慮人民另有資訊來源同時監督政府與法院\pg{790 頁}；在這一設定下，監督者自身的代理成本無從被人民察覺，只能事前以獨立與否的選擇一次權衡。\footcite[788, 790]{law2009}
\end{enumerate}

\section{課前可提的問題}
\label{sec:70_independence-4}

\begin{enumerate}
\tightlist
\item 787 頁承認法院是追求自己目標的代理人，788 至 789 頁卻只把這項成本表現為判違憲的比率偏高。法院的偏向若是選擇性的，Law 的權衡是否仍然成立？
\item 790 頁所說「不想要的司法造法」由誰承擔？人民對法院誤判違憲的容忍，是否以人民偏好一致為前提？
\end{enumerate}
